\section{Historical Matched-Selection Construction}
\label{sec:method}

\paragraph{Historical scope.} This appendix restores the original matched-content design, equations, and algorithm. Its 32-cooling/32-heating panel and four/six-initialization studies are historical; they are not the main 64-placement/64-cleaning panel or its eight-initialization global comparison.
\begin{figure}[htbp]\centering
\resizebox{\linewidth}{!}{\begin{tikzpicture}[
  font=\sffamily\footnotesize,
  block/.style={draw=ink!55,rounded corners=2pt,align=center,text width=2.72cm,minimum height=1.05cm,inner sep=5pt,fill=white},
  flow/.style={-{Stealth[length=2mm]},line width=.65pt,draw=ink!75},
  control/.style={flow,dashed,draw=agentblue},
  note/.style={align=center,font=\sffamily\scriptsize,text=ink!80}]
\node[block,fill=ink!4] (harness) at (0,0) {Historical collector\\base agent + $H_s$};
\node[block,fill=ink!4] (env) at (3.6,0) {ALFWorld\\text environment};
\node[block,fill=ink!4] (pool) at (7.2,0) {Successful episodes\\fixed source pool $\D_s$};
\node[note,text width=2.7cm] at (10.8,0) {API and Generic:\\two source pools,\\one environment};
\draw[flow] (harness.north east) to[bend left=18] node[above,note]{action} (env.north west);
\draw[flow] (env.south west) to[bend left=18] node[below,note]{observation} (harness.south east);
\draw[flow] (env) -- (pool);
\node[block] (view) at (0,-2.0) {Student-view examples\\real history + action};
\node[block,fill=agentblue!6] (score) at (3.6,-2.0) {Frozen base student\\action NLL $q_0(u)$};
\node[block,fill=agentorange!7] (select) at (7.2,-2.0) {Matched selection\\random vs.\ low NLL};
\node[block,fill=agentblue!6] (train) at (10.8,-2.0) {Paired LoRA SFT\\same initialization $i$};
\draw[flow] (pool.south) -- ++(0,-.35) -| (view.north);
\draw[flow] (view) -- (score);
\draw[flow] (score) -- (select);
\draw[flow] (select) -- (train);
\node[block,text width=4.05cm,fill=agentblue!5] (quotas) at (1.8,-4.0) {Frozen reference schedule $\mathbf r_j$\\task/action positions + support quotas};
\node[block,text width=3.25cm] (result) at (6.55,-4.0) {Paired task success $\Delta_s$\\all cells + per-initialization means};
\node[block] (eval) at (10.8,-4.0) {Common executor\\assistance withdrawn};
\draw[control] (quotas.north east) -- ++(.35,.30) -| (select.south);
\draw[flow] (train) -- node[right,note]{256 distinct records\\768 presentations\\192 updates} (eval);
\draw[flow] (eval) -- (result);
\node[note,text width=12.8cm] at (5.4,-5.1) {Selection, training, and evaluation use no new teacher calls. Both policy arms are trained from the same base checkpoint.};
\end{tikzpicture}
}
\caption{The original historical pipeline, separating the collector, student-view examples, matched selection, paired training, and common post-training executor. It documents the predecessor experiment; the main global-selection architecture is Figure~\ref{fig:architecture}.}
\end{figure}

\subsection{Examples, predictability, and training schedules}
A fixed source pool $\D_s$ contains eligible action demonstrations $z=(x,h,a)$: task $x$, actual public interaction history $h$, and the next canonical action $a$. Here $a$ includes the supervised JSON serialization and its chat-template terminator; $c(a)$ parses the command's action family. Let $u(z)$ identify the complete student-view messages by exact content. Deduplication identifies literal message equality, not semantic equivalence. Every retained example preserves its real history; selecting examples does not synthesize a new interaction trajectory.

For the frozen base student $p_{\theta_0}$, action predictability is measured by token-average target NLL,
\begin{equation}
 \ell_0(z)=-\frac{1}{|a|}\sum_{v=1}^{|a|}
 \log p_{\theta_0}(a_v\mid h,a_{<v}).
 \label{eq:nll}
\end{equation}
The original cache can assign different numerical values to identical content in different batching layouts. The historical pre-canonical selection experiments used the recorded mean over copies of each content identity,
\begin{equation}
 q_0(u)=\frac{1}{|\{z\in\D_s:u(z)=u\}|}
 \sum_{z\in\D_s:u(z)=u}\ell_0^{\mathrm{cache}}(z).
 \label{eq:cache}
\end{equation}
We retain this definition for historical provenance. The completed uniform-scoring diagnostic and canonical retraining are reported in Appendix~\ref{sec:scoring}; the fresh-panel main comparisons use the canonical scores.

A sampling recipe $j$ defines a frozen reference schedule $\mathbf r_j=(r_1,\ldots,r_B)$ with $B=768$ presentations. For a stratum function $g$, define the number of presentations and distinct content identities assigned to stratum $g$:
\begin{equation}
 n_g=\sum_{k=1}^{B}\ind\{g(z_{r_k})=g\},\qquad
 K_g=\left|\{u(z_{r_k}):g(z_{r_k})=g\}\right|.
 \label{eq:quotas}
\end{equation}
The two policies select $K_g$ identities from the eligible distinct support $\U_g$. Random selection samples uniformly without replacement; low-NLL selection solves
\begin{equation}
 A_g^{\low}\in\underset{A\subseteq\U_g,\,|A|=K_g}{\arg\min}
 \sum_{u\in A}q_0(u).
 \label{eq:selection}
\end{equation}
Ties are resolved by the representative corpus index. Within each stratum, a fixed recipe-specific shuffle cycles through the selected identities, visiting every identity before repeating it. These stratum sequences are interleaved in reference order. For the resulting policy schedule $\boldsymbol\sigma^{p}_{s,j}$,
\begin{align}
 g(z_{\sigma_k^p})&=g(z_{r_k})\quad &&\text{for every presentation }k,\label{eq:positions}\\
 \left|\{u(z_{\sigma_k^p}):g(z_{\sigma_k^p})=g\}\right|&=K_g\quad &&\text{for every stratum }g.\label{eq:support}
\end{align}
The selected content may differ while these constraints hold exactly.

\subsection{Progressively stronger controls}
\begin{figure}[t]
\centering
\resizebox{\linewidth}{!}{\begin{tikzpicture}[font=\sffamily\footnotesize,
 card/.style={draw=ink!40,rounded corners=2pt,align=left,text width=3.85cm,minimum height=1.65cm,inner sep=7pt},
 flow/.style={-{Stealth[length=2mm]},draw=ink!65,line width=.65pt}]
\node[card,fill=ink!3] (s2) at (0,0) {\textbf{1. Fix training dose}\\Full pass vs.\ subset replay\\768 draws / 192 updates\\\textcolor{ink!65}{Task mixture can differ}};
\node[card,fill=agentblue!5] (s3) at (4.65,0) {\textbf{2. Match task exposure}\\Random vs.\ low NLL\\Same task at each position\\256 distinct records};
\node[card,fill=agentorange!6] (s4) at (9.3,0) {\textbf{3. Match task/action}\\Random vs.\ low NLL\\Same task/action positions\\Same stratum support};
\draw[flow] (s2) -- (s3);
\draw[flow] (s3) -- (s4);
\end{tikzpicture}
}
\caption{\textbf{Progressive controls, not interchangeable treatments.} Each comparison answers a different conditional question. The final constraint also restricts eligible content support and weakens the low-NLL selection contrast; changes in effect size across stages cannot be read as a mediation fraction.}
\label{fig:controls}
\end{figure}
\smallpara{Training dose: complete pass versus subset replay.}
Both arms receive 768 presentations and 192 updates. Complete-pass scheduling traverses the eligible corpus before repeating rows; the replay arm reuses a 256-row subset. Membership, multiplicity, order, and task mixture can change together. This first comparison evaluates the scheduling intervention as a whole.

\smallpara{Task-matched selection.}
We set $g(z)=x$. The policies share task identity at every presentation, task counts, and 256 distinct task/content records. They can still differ in action composition. We cross four initialization/dropout seeds with four sampling recipes within each source.

\smallpara{Task-by-action-matched selection.}
We set $g(z)=(x,c(a))$ and inherit the task/action quotas from the task-matched random reference. The policies now share every task/action position and within-stratum support count. Action families absent from the reference remain absent. The support restriction retains approximately 66.15\% and 40.44\% of the preceding NLL-selection contrast in the API and Generic pools. It changes the intervention as well as the control. We cross six initialization/dropout seeds with the same four sampling recipes.

\begin{algorithm}[t]
\caption{Matched content selection and paired evaluation}
\label{alg:protocol}
\begin{algorithmic}[1]
\Require Fixed pool $\D_s$, base student $\theta_0$, reference $\mathbf r_j$, strata $g$
\State Compute content identities, scores $q_0(u)$, and quotas $(n_g,K_g)$
\For{$p\in\{\rand,\low\}$}
  \State Select $A_g^p$ independently within each reference stratum
  \State Cycle shuffled identities within each stratum; interleave in reference order
  \State Verify Eqs.~\eqref{eq:positions}--\eqref{eq:support}, total support, and schedule hashes
  \For{initialization/dropout seed $i$}
    \State Train a fresh LoRA adapter on $\boldsymbol\sigma^p_{s,j}$ for 192 updates
    \State Audit weights and evaluate through the common executor without collection guidance
  \EndFor
\EndFor
\State Report paired differences, all initialization means, and all wins/ties/losses
\end{algorithmic}
\end{algorithm}

\subsection{Training and estimands}
We fine-tune low-rank adapters~\citep{hu2022lora} while freezing the backbone. The supervised objective averages the token-normalized target loss over scheduled examples,
\begin{equation}
 \mathcal L(\phi;\boldsymbol\sigma)
 =\frac1B\sum_{k=1}^{B}\ell_{\theta_0+\Delta(\phi)}(z_{\sigma_k}).
 \label{eq:sft}
\end{equation}
The implementation uses microbatches of one example and gradient accumulation to a global batch of four. Both policies use the same initialization/dropout seed $i$, optimizer configuration, and update budget. Equal presentations do not imply identical input/target token totals or exact FLOPs.

Let $Y_{s,p,i,j,f,t}$ denote success for source $s$, policy $p$, initialization $i$, recipe $j$, family $f$, and task $t$. Each model's outcome and the primary contrast are
\begin{align}
 R_{s,p,i,j}&=\frac12\sum_{f\in\{\mathrm{cool},\mathrm{heat}\}}
                  \frac1{32}\sum_{t=1}^{32}Y_{s,p,i,j,f,t},\label{eq:success}\\
 \Delta_s&=\frac1{IJ}\sum_{i=1}^{I}\sum_{j=1}^{J}
                (R_{s,\low,i,j}-R_{s,\rand,i,j}).\label{eq:effect}
\end{align}
Initialization seeds and sampling recipes are crossed factors. We also report $\Delta_{s,i}=J^{-1}\sum_j(R_{s,\low,i,j}-R_{s,\rand,i,j})$. The 24 task/action-matched cells are not 24 independent initialization seeds. Conditional repeats are excluded from these primary means.
